GRPO-trained neurosymbolic formalizers routinely achieve high prove rates on logical reasoning benchmarks, yet no principled diagnostic exists to assess whether the underlying formal outputs contain genuine logical structure or degenerate shortcuts.
We introduce \emph{Solver Trace Ablation} (STA), a solver-agnostic protocol that derives the \emph{Intensional Structure Ratio} (ISR) via leave-one-out component necessity testing: each premise is ablated and the solver re-executed, yielding a continuous score measuring the fraction of components genuinely required for the proof.
Applied to GRPO training of a Qwen3-4B formalizer across 6 reward configurations, 3 random seeds, and up to 1{,}000 training steps, STA reveals three findings.
First, ISR serves as an effective diagnostic: execution-only GRPO achieves 96\% prove rate while ISR collapses from 0.64 to 0.003 over training---a 99.5\% decline quantifying structural degeneration invisible to the execution reward.
All three seeds exhibit 87--100\% conclusion embedding in proved outputs; a supervised fine-tuning baseline confirms that these pathologies are RL-specific (92.9\% genuine proofs, near-zero hacking).
Second, ISR fails as a training signal due to meta-level Goodhart's Law: at low dose ($\beta{=}0.1$), the policy games ISR itself via single-premise conclusion duplication (64.5\% H7 across 3 seeds vs.\ 0\% under matched random controls), with the apparent structural advantage dissolving entirely under premise-count-matched analysis (1.06$\times$); at high dose ($\beta{=}0.3$), classic Goodhart collapse drives prove rate to 18\% versus 48\% for the random control.
Third, both the diagnostic capability and the Goodhart failure transfer to Prolog (SWI-Prolog), where ISR-as-reward produces \emph{universal} Goodhart: all 3 seeds converge to gaming with 0\% genuine output, each via a distinct mechanism---tautological duplication, degenerate syntax, and query-as-fact gaming.
Across all conditions we identify and quantify 7 distinct hacking modes, providing an empirically grounded taxonomy for structural failures in neurosymbolic RL.
Our results expose a fundamental asymmetry: ISR is reliable as a diagnostic observer but unreliable as an optimization target---a cautionary finding for any proxy-metric-as-reward approach.
